
\subsection{Exact resources of the prescribed expansion}

\begin{remark}[Reading the resource bounds]
The content of Theorem~\ref{thm:resA-resources} for most readers is (RA1),
(RA2) and (RA4): in the expansion that literally writes body, keep,
reversed body at every mark, work doubles per nesting level, exactly;
retained wires do not multiply, because inverses reuse the same named
records. Two warnings govern everything in this section. First, ``tight''
always means tight for that prescribed, unsimplified expansion: with an
empty keep the expanded circuit can already denote the identity, so
nothing here lower-bounds the best equivalent circuit, parallel depth,
or time at bounded space. Second, the remaining clause, (RA3), is a
routine but lengthy verification that a classical Turing machine can
print the expansion without random access; it is logically separate from
(RA1)--(RA2) and can be skipped on a first reading. No pebble-game
trade-off continuum is proved in this paper: the only proved alternatives
to eager expansion are the restricted delayed construction of
Theorem~\ref{thm:resA-delayed} and deletion of an inner inverse exactly
when the commutator test (RA9) (or its subspace form (RA10)) holds.
\end{remark}

\begin{theorem}[Expansion resources, not optimal synthesis]
\label{thm:resA-resources}
Assume all hypotheses of Theorem~\ref{thm:sA-sound}, including finite
static syntax, total completions, retained logical registers, the
exclusions of Definition~\ref{def:tiera}, and specified exact clean
circuits over $\mathcal G^{\pm}$. Use Definition~\ref{def:costs}, with
no cancellation or circuit simplification. In particular all tick
decoders satisfy the canonical-field, live-flag and fault requirements
on every code, and exclude allocation and measurement modes on every
code. Marks pair statically; marked bodies contain no allocation, drop,
observed measurement, external console observation or actual stop.
The menu extension requires base Tier~A entries with supplied exact
finite-alphabet lists and retained disjoint selectors. A compiled menu
is one leaf for $C,L,d$ and for the following sum: its payload gate
costs are not counted again. For an atom occurrence $v$
let $h(v)$ count its enclosing marks. For a keep occurrence $v$ let
$h(v)$ count its strictly enclosing marks, excluding the mark whose
keep it is. Write $c_v$ for its supplied circuit length. Then
\[
 \begin{gathered}
 C(I)=0,\qquad C(a)=c_a,\qquad C(P;R)=C(P)+C(R),\\
 C(\mathsf M(P,K))=2C(P)+c_K,\qquad
 C(P)=\sum_v2^{h(v)}c_v\leq2^{d(P)}L(P).             \tag{RA1}
 \end{gathered}
\]
The sum is an exact count, hence both an upper and a matching lower
bound for this expansion. Under serial execution, time in gate ticks is
$C(P)$ by the definition of that unit. Define $q_{\rm in}$ as the
number of data input wires in the enclosing outer program, so that its
fixed data pool has $q=q_{\rm in}+q_{\rm new}(P)$ wires.
Classical preparation of the supplied gate lists is a separate cost;
every gate in those lists is included in the execution count.

For the snapshot representation assume one shared $s$-bit configuration,
common sample width $k\geq1$, and common upper bounds $\beta,\Delta$
on the device circuits of Definition~\ref{def:device} across all tick
occurrences. Use globally distinct syntactic records. A sufficient
total wire count is
\[
 \begin{split}
 N(P)={}&q_{\rm in}+q_{\rm new}(P)+s+(s+k)n_{\rm tick}(P)\\
       &{}+s n_{\rm cont}(P)+m_{\rm extra}
         +\sum_\nu\delta_\nu+a_* .
 \end{split}                                         \tag{RA2}
\]
Here the sum counts disjoint retained menu selectors of widths
$\delta_\nu$, excluded from $m_{\rm extra}$ by convention in this
display. Every payload tick and continuation occurrence is counted once
for record allocation, across all entries and menu occurrences. If a retained
selector is computed by a subterm of $P$, that subterm's tick and continuation
records are already included in $n_{\rm tick}(P)$ and $n_{\rm cont}(P)$,
any other logical wires it retains are included in $m_{\rm extra}$,
and its clean scratch is included in $a_{\rm other}$ (or in the tick
and keep terms of (RA2a) when it consists of those constructions).
A selector supplied from outside $P$ contributes only its retained width
$\delta_\nu$ here; the cost of computing it is not part of this display. Inverses reuse these names.
The maximum local clean scratch requirement is $a_*$. For the concrete
implementations constructed below it can be chosen to satisfy
\[
 a_*\leq\max\left\{
 a_{\rm other},\ c_0(\beta+\Delta+s+k+1),\
 \max_K c_0\kappa_K,\
 \max_\nu\bigl(a_\nu^{\rm pay}+c_0(\delta_\nu+1)\bigr)
 \right\}.                                         \tag{RA2a}
\]
The absolute constant $c_0$ is sufficiently large. For a keep $K$,
$\kappa_K$ is the maximum, over its Boolean constituent circuits, of
one plus gate count plus input and output widths; an empty maximum
is zero. The number $a_\nu^{\rm pay}$ is the maximum clean workspace
of the supplied entry lists. Its base Tier~A constituents obey the
same bounds without the menu term. The declared number $a_{\rm other}$
bounds scratch for all other supplied atoms, including completed RTM
steps; no complexity bound for arbitrary supplied completions is assumed.
If arbitrary supplied circuits are used instead of these constructions,
their declared scratch widths must replace the corresponding bounds.
If every non-input wire in (RA2) is freshly zero-initialized and each
initialization costs one tick, add $N(P)-q_{\rm in}$ ticks; preparation
of a nonzero initial machine configuration is additional.
For
\[
 \mathcal R=\beta+\Delta+s+k+(q^2+kq+1)(s+k+1),
\]
a supplied fragment of $T_0$ joint ticks, including any declared timeout
tick, has $T_0(s+k)$
record bits and an implementation of $O(T_0\mathcal R)$ gates, before
multiplication by enclosing-mark factors in (RA1). The claim of ordinary
behavior at each forward tick additionally requires its logical records
blank at that use. A static menu with $0\leq j\leq2^\delta$ entries
of supplied exact clean $\mathcal G^{\pm}$ lengths $g_i$ and
a preserved $\delta$-bit selector disjoint from all payload and scratch
wires costs
$O(\sum_{i<j}g_i+j(\delta+1))$ gates, with clean scratch for one
selector flag and its computation of width $O(\delta+1)$, in addition
to the maximum payload scratch used simultaneously with that flag.

Use an explicit uncompressed ordered syntax tree encoded in prefix order
by constructor tokens with their arities and delimited atom/keep lists. Gate descriptions
are nonempty finite strings with a constant-width gate tag and binary
wire labels. If $\sigma$ is the total bit length of this input and
$\ell$ the maximum gate-description length (zero if there are no gates),
a multitape classical TM can print the
expansion in
\[
 O((\sigma+C(P)+1)^3(\ell+1)^2)                       \tag{RA3}
\]
steps. This deliberately loose bound excludes construction of the supplied
lists, and does not assume unit-cost random access.
For $P_0=a$, $c_a>0$, and $P_{j+1}=\mathsf M(P_j,K)$ with a fixed
keep circuit of length $c_K\geq0$,
\[
 C(P_j)=2^jc_a+(2^j-1)c_K.                            \tag{RA4}
\]
Thus the exponential dependence on nesting depth is attained for the
specified expansion. These statements give no lower bound on the
smallest equivalent circuit, circuit depth under parallel execution,
or optimal time under a space restriction.
\end{theorem}
\begin{proof}
 \step{resA-count-ind}{The recurrences and weighted sum (RA1) hold.}
 \begin{proof}
  \step{resA-count-base}{\case{Identity or an atom.}}
  \begin{proof}\pf\ Identity prints no gates. An atom prints its
  supplied list once, and has $h=0$.\end{proof}
  \step{resA-count-seq}{\case{A sequence, with the assertion for its
  two proper subterms.}}
  \begin{proof}\pf\ Concatenation adds lengths and retains each
  occurrence's number of enclosing marks.\end{proof}
  \step{resA-count-mark}{\case{$\mathsf M(P,K)$, with the assertion
  for $P$.}}
  \begin{proof}\pf\ The list is the body, keep, and reversed adjoint
  body. The alphabet is closed under adjoints, so reversal preserves
  length. Every body occurrence has one additional enclosing mark and
  appears twice as often; the new keep appears once with $h=0$.
  This proves both the recurrence and the weighted sum.\end{proof}
  \qedstep
  \begin{proof}\pf\ Structural induction gives the equality. For
  every occurrence $h(v)\leq d(P)$ and $c_v\geq0$, so summing
  $2^{h(v)}c_v\leq2^{d(P)}c_v$ gives the inequality.\end{proof}
 \end{proof}
 \step{resA-count-wires}{Equation (RA2) suffices without assuming
 logical machine cleanliness.}
\begin{proof}
 \step{resA-count-wires-logical}{The logical-wire contribution in
 (RA2) counts each required retained factor once.}
 \begin{proof}\pf\ Count the data wires, shared $s$-bit
 configuration, one $s+k$-bit cell per joint-tick occurrence, one
 $s$-bit cell per continuation occurrence and declared extra logical
 wires. An occurrence's adjoint reuses its names, so repetitions
 in (RA1) add no wires. All logical wires remain present.\end{proof}
 \step{resA-count-wires-selectors}{Menu selectors and payload records
 have exactly the declared additional allocation.}
 \begin{proof}\pf\ Add the disjoint selector widths explicitly.
 Selector-computation records and other logical wires are already
 included in the declared counts. Although a menu gate list is
 one leaf, its payload record declarations remain logical wires
 and are counted as stipulated in (RA2).\end{proof}
 \step{resA-count-wires-scratch}{One implementation pool of width
 $a_*$ suffices independently of logical-record reset.}
 \begin{proof}\pf\ Every atom/keep implementation and its inverse
 restores implementation scratch on every logical input.
 Lemma~\ref{lem:scratch} therefore permits sequential reuse of
 a pool of the maximum width. This assumption concerns
 implementation scratch, not cleanliness of logical machine records.
 \end{proof}
 \step{resA-count-wires-init}{The optional zero-initialization charge
 is $N(P)-q_{\rm in}$.}
 \begin{proof}\pf\ These are exactly the non-input wires in (RA2),
 each charged one tick under the stated convention. Preparation of
 a nonzero machine configuration is additional.\end{proof}
 \qedstep
 \begin{proof}\pf\ Add \stepref{resA-count-wires-logical}--\stepref{resA-count-wires-scratch}
 for the declared allocation. The following construction bounds the
 scratch width; \stepref{resA-count-wires-init} gives the optional
 initialization cost.\end{proof}
\end{proof}
 \step{resA-count-tick}{The asserted tick and menu bounds hold.}
 \begin{proof}
  \step{resA-count-ticklocal}{One joint tick costs $O(\mathcal R)$.}
  \begin{proof}
   \pf\ We reuse the circuit of Lemma~\ref{eqv:tick-cost}, here with the
   sharper simultaneous-scratch count needed for (RA2a).
   \step{resA-count-decoder}{For a device tick, its quantum action
 $D_t$ has gate cost
 $O(\Delta+(q^2+kq+1)(s+k+1))$ and scratch
 $O(\Delta+s+k+1)$.}
\begin{proof}
 \step{resA-count-decoder-eval}{Computing the decoder costs
 $O(\Delta)$ gates and workspace.}
 \begin{proof}\pf\ Apply Lemma~\ref{lem:synthesis} to the canonical
 decoder on unchanged $C$; $\Delta$ includes its input and output
 ports. No global validity test is added: canonical output, legal
 live operands and fault handling on every code are hypotheses.
 Measurement and allocation modes are excluded here.\end{proof}
 \step{resA-count-decoder-predicates}{There are $O(q^2+kq+1)$
 predicates, each of size and workspace $O(s+k+1)$.}
 \begin{proof}
  \step{resA-count-decoder-gatepred}{There are $O(q^2+1)$ gate
  predicates, each of Boolean size $O(a_q+1)$.}
  \begin{proof}\pf\ Enumerate $O(q^2+1)$ gate actions. A predicate
  compares only mode, gate tag and at most two addresses. Equality
  is a conjunction of bit comparisons, of size $O(a_q+1)$.\end{proof}
  \step{resA-count-decoder-readpred}{There are $kq$ read
  predicates, each of Boolean size $O(a_q+a_k+1)$.}
  \begin{proof}\pf\ Enumerate $kq$ position/address pairs. Test mode,
  the one read address and whether the fixed position is below the
  read length. Scan the length from most to least significant bit,
  retaining equality and less-than bits. This has Boolean size
  $O(a_q+a_k+1)$.\end{proof}
  \qedstep
  \begin{proof}\pf\ Definition~\ref{def:device} gives $a_q\leq s$,
  and $a_k=\lceil\log_2(k+1)\rceil\leq k$ for $k\geq1$.
  Lemma~\ref{lem:synthesis} gives the same order of workspace,
  independently of the entire decoder's width.\end{proof}
 \end{proof}
 \step{resA-count-decoder-apply}{Each selected action has constant
 additional gate cost and scratch.}
 \begin{proof}\pf\ Hold its predicate flag and apply controlled-$H$,
 controlled-$T$ or controlled-CNOT by Lemma~\ref{lem:synthesis}.
 Targets lie in $QB_t$, disjoint from $C$, decoder, flag and evaluator
 workspace. Gate mode selects one gate; read mode selects its listed
 copies, whose distinct sample targets make their order immaterial.
 \end{proof}
 \step{resA-count-decoder-erase}{One predicate computation and
 constant gate scratch suffice alongside the decoder word.}
 \begin{proof}\pf\ Erase each predicate before reusing its workspace;
 its decoder inputs have not changed. After all actions erase the
 decoder, whose input $C$ has not changed. Each erasure is the
 inverse evaluation and has the same cost.\end{proof}
 \qedstep
 \begin{proof}\pf\ Sum the gate costs over
 $O(q^2+kq+1)$ predicates and add decoder evaluation and erasure.
 Sequential reuse in \stepref{resA-count-decoder-erase} gives
 scratch $O(\Delta+s+k+1)$.\end{proof}
\end{proof}
   \step{resA-count-update}{The update and the timeout obey the
   remaining asserted bounds.}
   \begin{proof}\pf\ After $D_t$, evaluate $f(C,B_t)$ into $J_t$,
   undo its Boolean workspace, then swap $C,J_t$. By
   Lemma~\ref{lem:synthesis} this costs $O(\beta+s+k)$ gates and
   scratch; a swap is three CNOTs per bit. It is the required $F_t$
   even on nonblank logical records. The separate timeout has
   $D_t=\id$ and update $h_\tau$: equality of the status with its
   active code, followed by conditional constant replacement, has
   $O(s)$ Boolean size. The same snapshot construction costs
   $O(s+k)$ and keeps its declared $s+k$ recording bits. A continuation
   replaces a fixed position field and costs $O(s)$ gates and scratch.
   Each is bounded by the tick term in (RA2a).\end{proof}
   \qedstep
   \begin{proof}\pf\ By \stepref{resA-count-decoder} and
   \stepref{resA-count-update}, sequential reuse gives $O(\mathcal R)$
   gates and $O(\beta+\Delta+s+k+1)$ scratch per tick.
   The $T_0$ listed ticks, including any timeout, have distinct cells.
   Their lengths add; blankness for forward interpretation is exactly
   the promise in Definition~\ref{def:joint-tick}.\end{proof}
  \end{proof}
  \step{resA-count-menu}{A menu entry costs $O(g_i+\delta+1)$.}
\begin{proof}
 \step{resA-count-menu-flag}{Compute equality of the selector with
 $i$ in one flag using $O(\delta+1)$ gates and workspace.}
 \begin{proof}\pf\ Conjoin the bit equalities and apply the clean
 Boolean evaluator of Lemma~\ref{lem:synthesis}. The capacity
 $j\leq2^\delta$ makes every $i<j$ a selector value, including
 the zero-width selector when it has an entry.\end{proof}
 \step{resA-count-menu-control}{Hold the flag while controlling the
 $g_i$ elementary gates, at constant overhead per gate.}
 \begin{proof}\pf\ Each supplied gate is $H,T,T^\dagger$ or CNOT.
 Lemma~\ref{lem:synthesis} supplies its exact controlled version.
 Selector and predicate workspace are disjoint from all targets,
 including payload implementation wires. Thus simultaneous
 workspace is $a_\nu^{\rm pay}+O(\delta+1)$, including the
 constant scratch of the controlled elementary gate.\end{proof}
 \step{resA-count-menu-false}{\assume{The flag is zero.}
 \prove{The complete controlled list is identity in this sector.}}
 \begin{proof}\pf\ Every controlled gate is identity in this
 sector; consequently the complete list is identity.\end{proof}
 \step{resA-count-menu-true}{\assume{The flag is one.}
 \prove{The complete controlled list equals the supplied payload
 circuit, including its phase, and restores its implementation scratch.}}
 \begin{proof}\pf\ The controlled list is exactly the supplied
 clean payload circuit, including its phase. Its implementation
 scratch is zero after the complete list.\end{proof}
 \step{resA-count-menu-reuse}{Erase the flag and reuse the scratch
 for the next entry.}
 \begin{proof}\pf\ The selector has not changed, so inverse
 evaluation erases the flag. The two sectors above restore payload
 scratch; linearity gives the same restoration on arbitrary selector
 superpositions. Lemma~\ref{lem:scratch} permits sequential reuse.
 \end{proof}
 \qedstep
 \begin{proof}\pf\ Add flag computation and erasure to the $g_i$
 controlled gates, then sum over entries for the asserted menu bound.
 Only supplied exact finite lists are used, so this covers no
 ideal-angle payload lacking such a list.\end{proof}
\end{proof}
  \step{resA-count-keep}{The remaining scratch bounds in (RA2a) hold.}
  \begin{proof}\pf\ For each Boolean constituent of a keep use the
  evaluator of Lemma~\ref{lem:synthesis} directly into the disjoint
  XOR target and erase its stored Boolean gate values. Its cost and
  scratch are $O(\kappa_K)$. Constituent keeps execute sequentially,
  so their maximum scratch suffices. Arbitrary other atom scratch is
  the supplied bound $a_{\rm other}$. Menu entries are base terms;
  Lemma~\ref{lem:scratch} bounds their scratch by the maximum of
  their base constituents. These observations and the preceding two
  steps establish (RA2a) for one fixed sufficiently large $c_0$.
  \end{proof}
  \qedstep
  \begin{proof}\pf\ By \stepref{resA-count-ticklocal} and
  \stepref{resA-count-menu} and \stepref{resA-count-keep}.\end{proof}
 \end{proof}
 \step{resA-count-compile}{A sequential-tape compiler has bound (RA3).}
 \begin{proof}
  \step{resA-compile-encoding}{\pick\ $N$ syntax-node identifiers
 of common width $N+1$, a child-before-parent order, and fixed-width
 gate fields of width $w=O(\ell+1)$. They can be prepared in
 $O((\sigma+1)^3)$ multitape steps, with $N=O(\sigma+1)$.}
\begin{proof}
 \step{resA-compile-encoding-ids}{There are $N=O(\sigma+1)$ nodes;
 assign the $i$th constructor the identifier $1^i0^{N+1-i}$.}
 \begin{proof}\pf\ Count constructor tokens in one input scan.
 Each consumes input symbols. The displayed distinct unary strings
 all have width $N+1$, and writing all of them costs $O((N+1)^2)$.
 \end{proof}
 \step{resA-compile-encoding-order}{A stack scan emits each node
 after its children, in $O((\sigma+1)^2)$ steps.}
 \begin{proof}
  \step{resA-compile-encoding-frames}{Use fixed-width $O(N+1)$ frames
  holding the identifier, constructor tag, up to two child identifiers
  and the number of children still awaited.}
  \begin{proof}\pf\ The prefix encoding supplies each constructor's
  arity. All identifier fields have width $N+1$, and the remaining
  fields have bounded width. Keep the stack head at its top.\end{proof}
  \step{resA-compile-encoding-leaf}{\assume{A leaf is completed.}
  \prove{Its record can be emitted, its frame popped, and its
  identifier recorded in the current parent frame, if any.}}
  \begin{proof}\pf\ It awaits no child. Emit its record, pop its
  frame and record its identifier in the current parent frame, if
  there is one.\end{proof}
  \step{resA-compile-encoding-child}{\assume{A nonleaf receives a
  completed child's identifier.}\prove{The child can be recorded;
  on the last child the node can be emitted and popped after all
  its children, and its identifier passed to its parent.}}
  \begin{proof}\pf\ Store the identifier and decrease the number
  awaited. On the last child, emit and pop the parent frame and
  pass its identifier to its parent. Thus all children have already
  been emitted when a node is emitted.\end{proof}
  \qedstep
  \begin{proof}\pf\ Each frame is pushed and popped once; there
  are $O(N)$ events. Reading or updating its constant number of
  fields, copying an identifier and returning the auxiliary tape
  head cost $O(N+1)$ per event. Add the input scan.\end{proof}
 \end{proof}
 \step{resA-compile-encoding-lists}{The owner-indexed tape of padded
 supplied lists has length $O((\sigma+1)^2)$.}
 \begin{proof}\pf\ Copy each delimited list under its owner
 identifier when encountered. Their original total length and the
 number of gates are at most $\sigma$, since descriptions are
 nonempty. A scan finds $\ell$; right-padding and copying cost
 $O((\sigma+1)(\ell+1))$. Owner headers cost $O((N+1)^2)$.
 Since $N,\ell\leq O(\sigma+1)$, the stated length follows.
 Fixed delimiters distinguish padded field ends.\end{proof}
 \step{resA-compile-encoding-lookup}{Retrieving a node's supplied
 list costs $O((\sigma+1)^2)$; all such retrievals cost
 $O((\sigma+1)^3)$.}
 \begin{proof}\pf\ Rewind and scan the owner-indexed tape.
 Compare unary fields synchronously and return the query head in
 $O(N+1)$ per header. This is proportional to header length, so
 it is included in the complete scan. Sum the tape-length bound
 over the $N$ nodes; charge each retrieval to its node below.\end{proof}
 \step{resA-compile-encoding-adjoint}{Adjoint conversion changes
 only a constant number of tag symbols in a gate field.}
 \begin{proof}\pf\ The constant-width tag ranges over
 $\{H,T,T^\dagger,\mathrm{CNOT}\}$. Its adjoint is a fixed table
 in finite control. Wire labels are unchanged, including the two
 labels of self-adjoint CNOT.\end{proof}
 \qedstep
 \begin{proof}\pf\ The preceding scans use a fixed finite number
 of tapes and no indexed access. They prepare the stated choices
 within $O((\sigma+1)^3)$ steps, with the subsequent list retrievals
 bounded separately by \stepref{resA-compile-encoding-lookup}.\end{proof}
\end{proof}
  \step{resA-compile-size}{For every processed node $R$,
  $C(R)\leq C(P)$; the expansion table has length at most
  $T=O((\sigma+1)(C(P)(\ell+1)+\sigma+1))$.}
  \begin{proof}\pf\ Every child length is at most its parent's
  length by (RA1) and nonnegativity. Following the finite ancestor
  chain proves the first assertion. Encode an expansion-table entry
  as its width-$(N+1)$ identifier, a delimiter, its expanded gate
  fields, and an end delimiter; table delimiters are distinct from
  gate symbols. Each entry has length $O(N+1+C(P)w)$. There are
  at most $N$ entries, giving $T$. Empty lists still have a header
  and an end delimiter.\end{proof}
  \step{resA-compile-node}{One node's expanded list is constructed
  and appended in $O(T+(\sigma+1)^2+(C(P)+1)(\ell+1)^2+1)$ steps.}
  \begin{proof}
   \step{resA-compile-lookup}{Every child lookup costs $O(T)$ steps.}
   \begin{proof}\pf\ Rewind the table tape and scan it sequentially.
   At each header compare its identifier against the query tape, then
   rewind the query head. This takes $O(N+1)$ per header. A failed
   match skips the entry, scanning each field once; a successful match
   copies the list to a separate tape as it is scanned. The header
   comparison and head return cost is bounded by a constant times
   header length. Thus a complete lookup, including the initial rewind,
   costs $O(T)$, not $O(NT)$. There are at most two child lookups
   per node.\end{proof}
   \step{resA-compile-reverse}{Reversed adjoint copying of an $m$-gate
   list costs $O((m+1)w^2)$.}
   \begin{proof}\pf\ First scan to the end, costing $O((m+1)w)$.
   From the end of each fixed-width gate field move left to its start,
   copy its symbols left-to-right to the output while changing its tag
   to the adjoint tag, then move left by at most $2w+O(1)$ to the
   preceding field's start. A fixed-width position counter may cost
   $O(w)$ per copied symbol even under a unary representation; resetting
   it costs $O(w)$ per field. This gives the stated conservative
   $O((m+1)w^2)$ bound including head movements. Each gate is visited
   once, in reverse gate order, while its wire-label bits retain their
   order. The zero-gate case only reads the delimiters.\end{proof}
   \qedstep
   \begin{proof}\pf\ Identity emits an empty list; an atom copies
   its supplied list; a sequence concatenates its two retrieved lists;
   a mark concatenates its body, supplied keep, and the reversed
   adjoint body. Locating the node's supplied list on the owner-indexed
   tape costs $O((\sigma+1)^2)$ by \stepref{resA-compile-encoding}.
   Only a constant number of tapes are needed because
   this is one node at a time, after its children have been stored.
   Every relevant supplied list and child list has at most $C(P)$
   gates. Sequential copying and the preceding reversal bound control
   all work except lookup. Rewinding the table and moving to its end
   to append the new entry takes $O(T)$ further steps. Apply
   \stepref{resA-compile-lookup} and \stepref{resA-compile-reverse}.
   \end{proof}
  \end{proof}
  \step{resA-compile-total}{The total bound is (RA3), including
  $C(P)=0$.}
  \begin{proof}\pf\ Process the prepared order once. If $t_i$ is
  the cumulative time after $i$ nodes, the preceding step gives
  $t_{i+1}-t_i=O(T+(\sigma+1)^2+(C(P)+1)(\ell+1)^2+1)$ uniformly in $i$.
  Sum over $N=O(\sigma+1)$ nodes and add preprocessing and the final
  root-list copy. The result is
  \[
   O\bigl((\sigma+1)T+(\sigma+1)^3+
       (\sigma+1)(C(P)+1)(\ell+1)^2\bigr).
  \]
  Substitute the bound on $T$ from \stepref{resA-compile-size};
  each summand is bounded by a constant times
  $(\sigma+C(P)+1)^3(\ell+1)^2$. When $C(P)=0$, there are still
  only $N$ node-processing events; empty body lists are stored once,
  not recursively regenerated for their two uses.\end{proof}
  \qedstep
  \begin{proof}\pf\ By \stepref{resA-compile-encoding}--\stepref{resA-compile-total}.
  \end{proof}
 \end{proof}
 \step{resA-count-tight}{Equation (RA4) and the claimed limitations hold.}
 \begin{proof}\pf\ The recurrence is $C_{j+1}=2C_j+c_K$ with
 $C_0=c_a$. Substitution and induction give (RA4). In particular,
 $C_j\geq2^jc_a$. Each counted gate must occur in an unchanged
 prescribed list, proving its matching lower bound and serial time.
 Printing requires at least one output action per gate because each
 gate description is nonempty. Definition~\ref{def:keep} permits a
 finite product with no entries: choose its empty circuit, so
 $K=\id$ and $c_K=0$. Already $P_1$ denotes
 $A_a^\dagger A_a=\id$, and every $P_j$, $j\geq1$, denotes
 identity. Its optimal gate count is zero although the expansion has
 $2^jc_a$ gates. Therefore an optimal-circuit lower bound does not
 follow. With $c_K=0$ the upper bound $2^dL$ itself is attained.
 \end{proof}
 \qedstep
 \begin{proof}\pf\ By \stepref{resA-count-ind}--\stepref{resA-count-tight}
 and the clean correctness of the expansion in
 Theorem~\ref{thm:sA-sound}.\end{proof}
\end{proof}

\subsection{Retaining intermediate computations}

\begin{theorem}[Time versus retained workspace]
\label{thm:resA-delayed}
Let $d\geq1$, let $b_0,\ldots,b_d,u_1,\ldots,u_d$ be nonnegative
integer widths, and let
$f_j:\{0,1\}^{b_{j-1}}\to\{0,1\}^{b_j}$ be total deterministic
functions. Let $X,R$ be disjoint external registers with
$|X|=b_0$, $|R|=b_d$. Supply total evaluator circuits $A_j$ of lengths
$a_j$ over $\mathcal G^{\pm}$, counted as in Definition~\ref{def:costs}.
On input $\ket x\ket{0^{b_j}}\ket{0^{u_j}}$ they produce
$\ket x\ket{f_j(x)}\ket{g_j(x)}$ for specified Boolean strings
$g_j(x)$ of width $u_j$. All evaluator implementation wires are included
in $u_j$; there is no additional implementation pool in this theorem.
For each $j$ supply a scratch-free $\mathcal G^{\pm}$ circuit of length
$k_j$ acting only on a $b_j$-bit output and a disjoint $b_j$-bit target,
with exact action $\ket{z,y}\mapsto\ket{z,y\oplus z}$ on every input.
The usual $b_j$ CNOTs are an admissible such list. Inverse lists have
the same length because this gate alphabet is adjoint-closed.
All evaluator output and private wires initially are zero.
Each evaluator acts only on its own input, output and private wires.
Only intermediate outputs, never the final external target, become
inputs of later evaluators. Earlier evaluator private wires remain
untouched until their specified inverse. No additional operation changes
retained values. Use the eager and delayed compositions of
Definition~\ref{def:delayed}, under these explicit hypotheses.
Let $F_j=f_j\circ\cdots\circ f_1$.

Both the eager and delayed constructions implement, on all states of
$XR$ with their specified auxiliary zero initialization, the same retained
isometry after padding their auxiliary lists to a common width and
permuting auxiliary factors to identify their names:
\[
 \ket{x,y}\ket0\longmapsto
 \ket{x,y\oplus F_d(x)}\ket0,                         \tag{RA5}
\]
including arbitrary external references. They need not be the same
unitary on nonzero auxiliary inputs. The exact unsimplified gate counts
and sufficient allocated auxiliary widths, excluding $X,R$, are
\[
 \begin{gathered}
 o_j=2a_j+k_j,\qquad E_1=o_1,\qquad E_j=2E_{j-1}+o_j,\\
 E_d=\sum_{j=1}^d2^{d-j}o_j,\qquad
 D_d=2\sum_{j=1}^da_j+k_d.                            \tag{RA6}
 \end{gathered}
\]
\[
 W_E=\sum_{j<d}b_j+\max_{1\leq j\leq d}(b_j+u_j),
 \qquad W_D=\sum_{j\leq d}(b_j+u_j).                 \tag{RA7}
\]
These are serial gate times as well as gate counts. If, uniformly in
$d,j$, $0<\alpha\leq a_j\leq A$ and $0\leq k_j\leq K$ for fixed
constants, then $E_d=\Theta(2^d)$ and $D_d=\Theta(d)$.
There is no universal strict separation between the two space counts;
they are sufficient allocations, not space lower bounds.

The same conclusions, including exact counts, hold for quantum
private outputs if the evaluator hypothesis is replaced by
\[
 A_j\bigl(\ket x\ket{0^{b_j}}\ket{0^{u_j}}\bigr)
   =\ket x\ket{f_j(x)}\ket{\gamma_{j,x}},             \tag{RA8}
\]
where each $\gamma_{j,x}$ is any unit vector, with the same widths,
alphabet, copying, locality, zero-initialization and no-intervention
assumptions. Thus quantum phases and non-basis private states are
permitted. The restrictions are exact preservation of the basis input,
an exact computational-basis output $f_j(x)$, and a product private
factor for each basis input. Disturbed inputs or an output entangled
with its private factor on a basis input are excluded.
\end{theorem}
\begin{proof}
 \step{resA-delay-oracle}{A compute--copy--inverse application of
 $A_j$ realizes a clean XOR by $f_j$ on its input and disjoint target.}
 \begin{proof}\pf\ On the basis input $\ket{x,y}$ and zero local
 wires, (RA8) gives
 $\ket x\ket{f_j(x)}\ket{\gamma_{j,x}}\ket y$.
 Copying changes only the last factor to $\ket{y\oplus f_j(x)}$.
 Applying $A_j^\dagger$ to the first three factors gives exactly
 $\ket x\ket0\ket0$, by total unitarity and (RA8).
 Linearity proves the assertion on arbitrary input/target states.
 The classical hypothesis is the special case
 $\gamma_{j,x}=\ket{g_j(x)}$.\end{proof}
 \step{resA-delay-eager}{The eager construction realizes (RA5).}
 \begin{proof}
  \step{resA-delay-eagerbase}{\case{$j=1$.}}
  \begin{proof}\pf\ Use \stepref{resA-delay-oracle}.\end{proof}
  \step{resA-delay-eagernext}{\assume{The previous composite is a
  clean XOR by $F_{j-1}$.}\prove{The level-$j$ composite is a clean
  XOR by $F_j$.}}
  \begin{proof}\pf\ Compute $F_{j-1}(x)$ into a zero intermediate.
  By the induction hypothesis every private wire used for that operation
  is reset. Apply the clean composite for $f_j$ with that intermediate as
  input and the current external target as target. It preserves the
  intermediate and changes the target by $f_j(F_{j-1}(x))$.
  Invert the previous composite to erase the intermediate; it acts
  independently of the changed target. This proves the basis action
  and zero restoration, then linearity proves the full isometry.\end{proof}
  \qedstep
  \begin{proof}\pf\ Induction using
  \stepref{resA-delay-eagerbase} and \stepref{resA-delay-eagernext}.
  \end{proof}
 \end{proof}
 \step{resA-delay-late}{The delayed construction realizes (RA5).}
 \begin{proof}\pf\ Fix $x$ and write $z_0=x$, $z_j=f_j(z_{j-1})$.
 After all forward evaluators, the retained vector is
 $\ket x\bigotimes_{j=1}^d
 (\ket{z_j}\ket{\gamma_{j,z_{j-1}}})$, with the target unchanged.
 This follows by induction from (RA8) and locality. Copy $z_d$ into
 the target. The factors of the last evaluator still have exactly
 its forward output, so its inverse erases them. Repeat in reverse
 order; every earlier input was preserved. The target is outside all
 these inverses. Hence all local wires are zero and (RA5) holds on
 each basis pair. Extend linearly and tensor with identity on a reference.
 No private vector is copied or discarded.\end{proof}
 \step{resA-delay-counts}{Equations (RA6)--(RA7) hold.}
\begin{proof}
 \step{resA-delay-counts-oracle}{The $j$th compute--copy--inverse composite has length
 $o_j=2a_j+k_j$.}
 \begin{proof}\pf\ Count its length-$a_j$ evaluator, length-$k_j$
 copy and length-$a_j$ inverse. Adjoint closure preserves list length.
 \end{proof}
 \step{resA-delay-counts-eager}{The eager counts are
 $E_1=o_1$, $E_j=2E_{j-1}+o_j$ and
 $E_d=\sum_{j=1}^d2^{d-j}o_j$.}
 \begin{proof}\pf\ The eager procedure uses the preceding
 composite twice and the current compute--copy--inverse composite once. The base is one
 oracle. Substituting the sum for $E_{j-1}$ into the recurrence
 doubles its coefficients and appends $o_j$, proving the formula
 by induction.\end{proof}
 \step{resA-delay-counts-delayed}{The delayed count is
 $D_d=2\sum_{j=1}^da_j+k_d$.}
 \begin{proof}\pf\ Each evaluator and its inverse occurs exactly
 once, and only the last copy occurs. Sum their list lengths.
 \end{proof}
 \step{resA-delay-counts-pool}{The eager procedure can share a local
 pool of width $\max_j(b_j+u_j)$.}
 \begin{proof}\pf\ Keep all intermediate outputs outside this
 pool. Step~\stepref{resA-delay-eager} restores the pool before
 the level-$j$ compute--copy--inverse composite; \stepref{resA-delay-oracle} restores it
 before the preceding composite's inverse. Lemma~\ref{lem:scratch}
 gives zero restoration on that inverse pass as well. Induction
 therefore justifies each reuse at every level. Pad unused pool
 wires by identity.\end{proof}
 \step{resA-delay-counts-widths}{The sufficient widths are
 $W_E=\sum_{j<d}b_j+\max_j(b_j+u_j)$ and
 $W_D=\sum_{j\leq d}(b_j+u_j)$.}
 \begin{proof}\pf\ For eager execution add the reserved intermediate
 wires to the shared pool of \stepref{resA-delay-counts-pool}.
 Delayed execution reserves a separate $b_j+u_j$ block for each
 evaluator. Neither allocation counts $X,R$. Renaming wires of
 a zero pool incurs no gates.\end{proof}
 \qedstep
 \begin{proof}\pf\ The first three steps give (RA6);
 \stepref{resA-delay-counts-widths} gives (RA7).\end{proof}
\end{proof}
 \step{resA-delay-orders}{The time orders and space qualifications hold.}
 \begin{proof}\pf\ The assumptions bound $o_j$ between $2\alpha$
 and $2A+K$. Sum $\sum_{j=1}^d2^{d-j}=2^d-1$ to obtain the eager
 order. Also $2\alpha d\leq D_d\leq2Ad+K$, proving the delayed
 order. For $b_j=b$, $u_j=u$, the displayed widths are
 $db+u$ and $d(b+u)$. They coincide when $u=0$ and differ by
 $(d-1)u$ otherwise. For general varying widths the sufficient
 formula $W_E$ need not even be less than $W_D$; no inequality
 between them has been used.\end{proof}
 \qedstep
 \begin{proof}\pf\ By \stepref{resA-delay-oracle}--\stepref{resA-delay-orders}.
 Classical basis private outputs satisfy (RA8), so this also proves
 the original restricted delayed-erasure target with its hypotheses.
 \end{proof}
\end{proof}

\subsection{The exact condition for deleting an inner inverse}

\begin{theorem}[A quantum deletion criterion and exact savings]
\label{thm:resA-deletion}
Let $A,P,B$ be specified total unitaries on a common finite retained
list $L=QM$, realized by Tier~A terms; let $K,K'$ be keeps of
Definition~\ref{def:keep}. There are no allocation, drop or observed
measurement operations in these terms. Consider the nested expression
\[
 \mathsf M\bigl(A;\mathsf M(P,K);B,K'\bigr).
\]
Here and below the same letters denote the terms' logical unitaries.
Delete the inner inverse $P^\dagger$ after $K$, retaining its records
until the outer inverse, to obtain $\mathsf M(A;P;K;B,K')$, where
$K$ is now a standalone atom. Write
\[
 U=BP^\dagger KPA,\quad V=BKPA,\quad H_0=B^\dagger K'B,
 \quad W=U^\dagger K'U,\quad \widetilde W=V^\dagger K'V.
\]
Then, with exact matrix phases,
\[
 W=\widetilde W\quad\Longleftrightarrow\quad
 [P,H_0]=0.                                          \tag{RA9}
\]
In particular it suffices that there is a tensor decomposition
$\cH_L=\cH_S\otimes\cH_{L\setminus S}$ over disjoint wire lists
with $P=P_S\otimes\id$ and $H_0=\id\otimes H_{L\setminus S}$,
even if $A,P,B$ contain quantum gates. More generally,
for any specified input isometry $J:\cH_{\rm in}\to\cH_L$,
equality of the retained isometries $WJ=\widetilde WJ$ is equivalent to
\[
 (PH_0P^\dagger-H_0)KPAJ=0.                          \tag{RA10}
\]
This weaker condition may suffice when (RA9) fails. These statements
are about retained operators first; any identical final partial trace
then preserves equality. They imply no reset of logical records on
their own and do not identify unequal dilations merely from equal
output channels.

For the resource assertion assume supplied exact clean circuits over
the adjoint-closed alphabet $\mathcal G^{\pm}$ as in
Definition~\ref{def:costs}. If $a,p,b,k,k'$ are the lengths of the
supplied expansions for $A,P,B,K,K'$, the two unsimplified counts are respectively
\[
 2(a+2p+k+b)+k',\qquad 2(a+p+k+b)+k',                 \tag{RA11}
\]
so precisely $2p$ gates and serial ticks are saved whenever the
transformation is justified. On a fixed retained $n$-wire logical list,
both can be implemented with $n+a_*^{\rm old}$ wires. Here
$a_*^{\rm old}$ is the original list's maximum local clean scratch;
the modified list uses only occurrences from that list, so
$a_*^{\rm new}\leq a_*^{\rm old}$.
Thus this local rule alone does not force a space increase. If it
replaces previously certified reuse of logical zero cells, those cells
must instead be retained separately or that reuse proved valid anew.
\end{theorem}
\begin{proof}
 \step{resA-delete-factor}{Put $F=KPA$. Then
 $W=F^\dagger PH_0P^\dagger F$ and
 $\widetilde W=F^\dagger H_0F$.}
 \begin{proof}\pf\ Substitute $U=BP^\dagger F$ and $V=BF$
 in the definitions and reverse factors under the adjoint.\end{proof}
 \step{resA-delete-full}{Equation (RA9) holds.}
 \begin{proof}\pf\ $F$ is unitary. Multiplying the two expressions
 in \stepref{resA-delete-factor} on the left by $F$ and right by
 $F^\dagger$ shows that equality is equivalent to
 $PH_0P^\dagger=H_0$. Multiplication on the right by $P$ makes
 this equivalent to $PH_0=H_0P$. Operators on disjoint tensor factors
 commute, proving the stated sufficient special case.\end{proof}
 \step{resA-delete-subspace}{Equation (RA10) is equivalent to
 $WJ=\widetilde WJ$.}
 \begin{proof}\pf\ Subtract the expressions in
 \stepref{resA-delete-factor}, multiply by $J$ on the right and
 by the invertible $F$ on the left. The result is exactly (RA10).
 The equality remains true when tensored with a reference identity;
 linearity and a final partial trace give the corresponding channel
 equality. Neither calculation asserts a zero record output.\end{proof}
 \step{resA-delete-resources}{The exact counts are (RA11).}
 \begin{proof}\pf\ The original outer body has length
 $a+2p+k+b$; the modified body has length $a+p+k+b$.
 Each is followed by its keep and its complete inverse, so apply
 (RA1). The logical names are unchanged; Lemma~\ref{lem:scratch}
 supplies the common clean scratch pool. Deletion removes occurrences
 both in the forward body and in its outer inverse; it introduces no
 new local implementation, proving the scratch inequality.
 Its hypothesis concerns
 implementation scratch, and does not certify logical cell reuse.
 \end{proof}
 \qedstep
 \begin{proof}\pf\ By \stepref{resA-delete-full}--\stepref{resA-delete-resources}.
 \end{proof}
\end{proof}

\begin{corollary}[Certified removal at multiple levels]
\label{cor:resA-flat}
Let a finite base Tier~A term, without derived menu atoms, have one
outermost mark with keep $K_*$. All resource lists are exact clean
$\mathcal G^{\pm}$ lists counted without simplification.
Let $S$ be the sum of supplied gate lengths of every atom and every
inner keep, each counted once. Suppose its expanded outer body can
be reduced to the list visiting each such occurrence once by successive
deletions of inner inverse lists. At each deletion, in the current
body write the unitary as $BP^\dagger F$, where the removed list
implements $P^\dagger$, and require
$[P,B^\dagger K_*B]=0$, or require (RA10) with $F$ in place of $KPA$
for the same fixed input embedding $J$. Retain every logical wire and
keep the supplied exact clean local circuits. Then the final outer
operation agrees with the original (on all retained inputs in the
commutator case, or on $J$ in the embedding case). Its count is
$2S+c_{K_*}$. Each deletion of a length-$p$ inverse list saves exactly
$2p$ gates. The time is therefore linear in this once-counted body size,
conditional on all the stated tests; no universal deletion rule is asserted.
\end{corollary}
\begin{proof}
 \step{resA-flat-one}{Each deletion preserves the specified equality.}
 \begin{proof}\pf\ The proof of Theorem~\ref{thm:resA-deletion}
 uses $F$ only as a unitary factor, so applies to the current factorization.
 Its two outer lists differ in length by twice the deleted length.
 \end{proof}
 \step{resA-flat-all}{The final count is $2S+c_{K_*}$.}
 \begin{proof}\pf\ Equality is transitive through the finite
 sequence of deletions. By hypothesis the final body contains each
 designated occurrence once, so has length $S$. Its outer expansion
 consists of this list, $K_*$, and its reversed adjoint list.\end{proof}
 \qedstep
 \begin{proof}\pf\ By \stepref{resA-flat-one} and \stepref{resA-flat-all}.
 \end{proof}
\end{proof}

\begin{counterexample}[A clean inner evaluation does not justify deletion]
\label{ex:resA-quantum}
Use five distinct qubits $q,r,t,a,b$, with $Q=qrt$, $M=ab$,
$A=\id$, and
\[
 P=\operatorname{CNOT}_{q\to a},\quad
 K=\operatorname{CNOT}_{a\to r},\quad
 B=\operatorname{CNOT}_{q\to b}H_q,\quad
 K'=\operatorname{CNOT}_{b\to t}.                     \tag{RA12}
\]
The Boolean copying gates here are explicit atoms, not full snapshot
translations of source reads. All atoms and keeps are total and allowed.
On the subspace supported on $a=0$, namely pure inputs
$\ket\psi_{qrtb}\otimes\ket0_a$ with factors placed by name and
their density-operator extensions, the inner marked operation
$P^\dagger KP$ is a clean CNOT from $q$ to $r$.
Nevertheless deleting its
inverse changes the outer retained isometry and even the reduced channel
on $qrt$. Here each channel uses the fixed isometry
$E:\cH_{qrt}\to\cH_{qrtab}$,
$E\psi=\psi_{qrt}\otimes\ket{00}_{ab}$ with factors placed by name,
then the respective retained operator, then $\tr_{ab}$; as channels
they are $\rho\mapsto\tr_{ab}(WE\rho E^\dagger W^\dagger)$ and
the same with $\widetilde W$. Let $\Pi_\pm=\proj\pm_t$ with
$\ket\pm=(\ket0\pm\ket1)/\sqrt2$. With identities and tensor
positions implicit, the full retained operators are
\[
 W=\Pi_++\Pi_-X_qX_rZ_b,\qquad
 \widetilde W=\Pi_++\Pi_-X_qX_aZ_b.                  \tag{RA13}
\]
For the input $\ket{00000}_{qrtab}$, their outputs, written in order
$q,r,t,a,b$, are respectively
\[
 \begin{split}
 W\ket{00000}&=
 (\ket{00}_{qr}\ket+_t+\ket{11}_{qr}\ket-_t)
               \ket{00}_{ab}/\sqrt2,\\
 \widetilde W\ket{00000}&=
 (\ket{00}_{qr}\ket+_t\ket0_a+
  \ket{10}_{qr}\ket-_t\ket1_a)\ket0_b/\sqrt2.
 \end{split}                                        \tag{RA14}
\]
In the first case the probability $r=1$ is $1/2$; in the second it is
zero, including after $M$ is traced. The prescribed counts are $11$
and $9$ gates, both on five logical wires with zero implementation
scratch. Omitting the Hadamard in $B$ makes the deletion valid on all
retained inputs. Thus its failure here depends on the intervening
quantum gate, rather than merely on nonzero initial auxiliary values.
\end{counterexample}
\begin{proof}
 \step{resA-ex-inner}{The inner marked operation is
 $R=P^\dagger KP=\operatorname{CNOT}_{a\to r}
                       \operatorname{CNOT}_{q\to r}$.}
 \begin{proof}\pf\ On a basis triple $(q,a,r)=(x,z,y)$, the three
 gates give $(x,z\oplus x,y)$, then
 $(x,z\oplus x,y\oplus z\oplus x)$, then
 $(x,z,y\oplus z\oplus x)$. This is the stated product.
 For $z=0$ it changes only $r$ by XOR with $x$ and restores $a$;
 linear extension gives cleanliness on the entire stated $a=0$
 subspace, including entanglement with untouched factors or a reference.
 \end{proof}
 \step{resA-ex-outer}{The conjugated outer keep is
 $H_0=B^\dagger K'B=\Pi_++\Pi_-X_qZ_b$.}
 \begin{proof}\pf\ $K'=\Pi_++\Pi_-Z_b$, since its target $X_t$
 has eigenvalues $+1,-1$. Conjugating $Z_b$ by
 $\operatorname{CNOT}_{q\to b}$ gives $Z_qZ_b$, as follows by
 multiplying a basis vector by $(-1)^{b\oplus q}$.
 Conjugation by $H_q$ changes $Z_q$ to $X_q$.\end{proof}
 \step{resA-ex-two}{Equation (RA13) holds.}
 \begin{proof}\pf\ The original operator is $R^\dagger H_0R$.
 In the expression for $R$ from \stepref{resA-ex-inner}, its
 $a\to r$ factor commutes with $X_q$, and conjugation by its
 $q\to r$ factor sends $X_q$ to $X_qX_r$ by direct basis action.
 For the modified body $V=BKP$, its outer operation is
 $P^\dagger K^\dagger H_0KP$. The keep $K$ commutes with $X_q$,
 whereas conjugation by $P$ sends $X_q$ to $X_qX_a$. All these
 gates commute with $Z_b$ and $\Pi_\pm$. Substitute into
 \stepref{resA-ex-outer} to obtain (RA13).\end{proof}
 \step{resA-ex-input}{Equation (RA14), the probabilities, and the
 failure of the commutator test follow.}
 \begin{proof}\pf\ Apply (RA13) to the all-zero basis vector, using
 $\Pi_\pm\ket0_t=\ket\pm_t/\sqrt2$ and $Z_b\ket0_b=\ket0_b$.
 The two displayed summands in either output are orthogonal and have
 squared norm $1/2$. In the original output only the second has
 $r=1$; in the modified output both have $r=0$.
 Moreover $PH_0P^\dagger=\Pi_++\Pi_-X_qX_aZ_b\ne H_0$,
 so (RA9) fails.\end{proof}
 \step{resA-ex-noH}{Without $H_q$ the deletion is valid, and the
 counts with $H_q$ are $11$ and $9$.}
 \begin{proof}\pf\ Without $H_q$, the conjugated keep is
 $\Pi_++\Pi_-Z_qZ_b$, which commutes with $P$ because CNOT
 preserves its control's $Z_q$ and does not touch $b,t$.
 Apply (RA9). With $H_q$, the lengths in (RA11) are
 $(a,p,b,k,k')=(0,1,2,1,1)$, giving $11$ and $9$.
 All five wires were explicitly listed and only elementary H and
 CNOT gates occur.\end{proof}
 \qedstep
 \begin{proof}\pf\ By \stepref{resA-ex-inner}--\stepref{resA-ex-noH}.
 \end{proof}
\end{proof}

\subsection{Small exact resource instances}

\begin{example}[Three depths with explicit circuit lists]
\label{ex:resA-counts}
In family N let $P_0=\operatorname{CNOT}_{q\to a}$,
$P_{j+1}=\mathsf M(P_j,\operatorname{CNOT}_{a\to r})$ on
$Q=qr$, $M=a$. Each supplied atom and keep has cost one; there
are three logical wires and no implementation scratch.
In family D let all $b_j=u_j=1$, $f_j(x)=x$, and let evaluator $A_j$
apply CNOT from its input to its zero output, then CNOT from that input
to its zero private bit. Thus $a_j=2$, $k_j=1$, $o_j=5$.
The external input and target are one qubit each. Workspace below means
the allocated sufficient widths (RA7), not an optimal width. All counts
include every gate of the prescribed lists, even cancellable pairs.
\[
\begin{array}{c|rrr|rrrr|rr}
 d& L_N&C_N&n_N&E_D&D_D&W_E&W_D&n_E&n_D\\\hline
 1&2&3 &3& 5 &5 &2&2&4&4\\
 2&3&7 &3&15 &9 &3&4&5&6\\
 3&4&15&3&35&13 &4&6&6&8
\end{array}                                         \tag{RA15}
\]
Here $n_N$ is total logical width of family N, and $n_E,n_D$ add the
two external qubits to family D's auxiliary widths. Serial gate time
is the corresponding gate-count column. Both family D procedures
realize the clean CNOT from their external input to their final target.
\end{example}
\begin{proof}
 \step{resA-table-N}{The three family N rows are correct.}
 \begin{proof}\pf\ There is one original atom and $d$ keeps, so
 $L_N=d+1$. Equation (RA4) with $c_a=c_K=1$ gives
 $C_N=2^{d+1}-1=3,7,15$. The gates use only $q,a,r$ and require
 no scratch; inverse repetitions do not add wires.\end{proof}
 \step{resA-table-D}{The three family D rows are correct.}
 \begin{proof}\pf\ On input $(x,0,0)$ each evaluator yields
 $(x,x,x)$, verifying its classical hypothesis and its cost two.
 A one-bit copy has cost one. Substitute into (RA6)--(RA7):
 $E_D=5(2^d-1)$, $D_D=4d+1$, $W_E=d+1$, $W_D=2d$.
 For $d=1,2,3$ these are exactly the displayed integers.
 Adding the external two wires gives $n_E=d+3$ and $n_D=2d+2$.
 The composite Boolean function is identity, so (RA5) is clean CNOT.
 \end{proof}
 \qedstep
 \begin{proof}\pf\ By \stepref{resA-table-N} and \stepref{resA-table-D}.
 \end{proof}
\end{proof}

\begin{remark}[Scope of the resource comparison]
The exponential lower count here is for a prescribed expansion, not for
quantum or reversible computation as a model. The delayed construction
has a proved retained semantics under its explicit hypotheses, including
its quantum extension (RA8). For arbitrary intervening gates, (RA9) or
(RA10) is a separate obligation at each proposed deletion; (RA12)
shows why cleanliness of an inner mark alone is insufficient. No
pebble-game theorem or unverified external reference carries a proof
step here. General optimal time--space trade-offs and unrestricted
dynamic mark control are not settled by these results.
The evaluator comparison requires exact basis outputs and a preserved
input even in (RA8), includes all evaluator wires in $u_j$, and uses
scratch-free XOR-copy lists. The multiple-deletion corollary is restricted
to base Tier~A terms without menus. The general resource theorem permits
menus only with supplied exact finite-alphabet lists and the once-only
gate/record accounting stated there. All restrictions collected in
Remark~\ref{rem:r3-restrictions} continue to apply.
\end{remark}
